\chapter{Current actions and derived interface centres}
\label{ch:vol2-currents}

A boundary action has to commute with fields inserted near the boundary.
When this commutation holds up to homotopy, two fields require a
compatibility between the two homotopies and their collision. The
compatibility is additional data. We first recover its cochains from
supported maps and construct them for a resolved current action.

The Heisenberg algebra makes the resulting distinctions explicit. At a
polynomial root, forgetting the current action retains many more
operators than the current cohomology sees. Their derived composition
has a nonzero Massey product. Over the level ring, the path in the
forgetful fibre turns a torsion current class into a free class. Finite
categories of interfaces then give centres in which cup products,
insertion brackets and parameter fibres can all be computed. These
calculations specify operations that a geometric bulk--boundary
comparison must preserve.

Throughout the chapter, $k$ is a field of characteristic zero. Complexes
are cohomological, with $K[1]^n=K^{n+1}$. For a homogeneous map $f$ the
Hom differential is $d_{\mathrm{Hom}}f=d f-(-1)^{|f|}f d$.
Polynomial states have finite support in their variables. A completion
will always be stated separately.

\section{Current homotopies from a supported map}
\label{sec:vol2-current-homotopies}

Let $O=k[y]$ and $D_y=O\langle\partial_y\rangle$, with
$[\partial_y,y]=1$. The Heisenberg vacuum at level $\kappa\in k^\times$
is
\begin{equation}\label{eq:vol2-fock-currents}
 V_\kappa=k[x_1,x_2,\ldots],\qquad
 J_{-n}=L_{x_n},\quad J_n=\kappa n\partial_{x_n}\quad(n>0),
 \qquad J_0=0.
\end{equation}
Here $L_f$ denotes multiplication by $f$. Its vacuum is $1$, its
current state is $J=x_1$, and its translation is
\[
 T=\sum_{n\geq1}n x_{n+1}\partial_{x_n},\qquad
 \mathcal V=O\otimes_kV_\kappa,\qquad \nabla=\partial_y+T.
\]
For a formal variable $u$, the field
$Y(J,u)=\sum_{m\in\mathbb Z}J_m u^{-m-1}$ determines the other
fields by normal products and translation. Direct calculation
gives $[J_m,J_n]=\kappa m\delta_{m+n,0}$ and
$[\nabla,J_m]=-mJ_{m-1}$. Each field is lower truncated on each state.

Take a second coordinate $z$ and put $u=z-y$. A map supported on $u=0$
takes values in principal parts, with all finite normal derivatives
retained as in Section~\ref{sec:vol2-principal-parts}. The current
generates a free $D_y$-submodule: $T^nJ=n!x_{n+1}$. Denote it by
$L_J=D_yJ$.

A semi-free complex has an increasing filtration by subcomplexes,
split on underlying graded modules, whose successive quotients
are direct sums of shifts of $D_y$.

\begin{lemma}[Laurent coefficients of supported maps]
\label{lem:vol2-laurent-current}
Let $Q$ be a semi-free complex of left $D_y$-modules. Let $j$ include
the complement of the diagonal and let $\Delta$ include the diagonal.
Then
\[
 R\operatorname{Hom}_{D_{z,y}}
   \bigl(j_*j^*(L_J\boxtimes Q),\Delta_+\mathcal V\bigr)
\]
is computed by ordinary differential-module maps. Taking the residue
identifies this complex with families of $O$-linear maps
$H_m:Q\to\mathcal V$, $m\in\mathbb Z$, satisfying
\begin{equation}\label{eq:vol2-laurent-horizontal}
 [\nabla,H_m]+mH_{m-1}=0.
\end{equation}
For each input, only finitely many $H_m$ with $m\geq0$ are nonzero.
There is no bound required uniformly over all inputs.
\end{lemma}
\begin{proof}
For one free generator the localized source has the telescope
presentation
\[
 D_{z,y}[u^{-1}]=\varinjlim_k D_{z,y}u^{-k},
\]
whose maps are right multiplication by $u$. The Ore localization
exists because $\operatorname{ad}u$ is locally nilpotent. Applying
Hom gives the inverse system of principal parts under multiplication
by $u$. These maps are onto: $u^{-r}$ has preimage $u^{-r-1}$.
Thus its ordinary limit computes its derived limit, by the product
model in the proof of Theorem~\ref{thm:vol2-continuous-chevalley}. Direct sums of
free generators become products of these exact calculations.

Use the semi-free filtration of $Q$, split on underlying graded
modules. At each finite stage the assertion follows by extensions.
The restriction maps between the ordinary Hom complexes are
degreewise onto, so the assertion passes to the entire filtration.
The resulting totalization uses products.

The map with residue family $(H_m)$ is reconstructed by
\begin{equation}\label{eq:vol2-laurent-reconstruction}
 \varphi(u^mJ\boxtimes q)
   =\sum_{r\geq0}H_{m+r}(q)u^{-r-1}.
\end{equation}
The stated finiteness makes this a principal part for each $m,q$.
Multiplication by $u$ shifts the formula. Linearity in $\partial_z$
then determines the values on every $T^nJ$. Finally
$\partial_y|_z=\partial_y|_u-\partial_u$ and
$\operatorname{Res}_u\partial_u=0$ give
\[
 \nabla H_m(q)=-mH_{m-1}(q)+H_m(\partial_yq),
\]
which is precisely the remaining differential-module relation.
These constructions are inverse and commute with the Hom differential.
\end{proof}

To extract these maps from a collision, retain its attaching map.
Write a one-current Cousin complex as $A[1]\oplus B$, where $A$ is
the open term, $B$ is the supported term, and $\mu:A\to B$ is the
unshifted residue map. For the suspension $s:A\to A[1]$ of degree
$-1$, its differential is
\[
 d(sa,b)=(-s\,d_Aa,\mu a+d_Bb).
\]
A map to a second such complex has open and closed components
$f_U,f_\Delta$ and a degree $-1$ cross component $h$. Its equation is
$d_{\mathrm{Hom}}h=f_\Delta\mu-\mu'f_U$.

Suppose that the prescribed unary operation is
$a:Q\to\mathcal V$, and that the open and closed restrictions are
identified with $1\boxtimes a$ and $a$ by specified homotopies
$\gamma,\delta$. Orient them by
$d\gamma=f_U-1\boxtimes a$ and $d\delta=f_\Delta-a$.
Writing $\partial_P,\partial_C$ for the source and target residue
maps, set
\begin{equation}\label{eq:vol2-corrected-current-face}
 \widetilde h=h+\partial_C\gamma-\delta\partial_P,
 \qquad
 d\widetilde h=a\partial_P-\partial_C(1\boxtimes a).
\end{equation}
For a semi-free $P$, take $Q=P\otimes_O\mathcal V$ with its
prescribed boundary current action. On a free $D_y$-summand,
the change from $\partial_y\otimes1$ to the diagonal connection
$\partial_y\otimes1+1\otimes\nabla$ is triangular in PBW degree.
It identifies the tensor with the module induced from the
$O$-free module $\mathcal V$, so the tensor is $D_y$-free.
Applying this to the filtration quotients makes $Q$ semi-free.
The residue map $\partial_P$
acts on the factor $\mathcal V$. The identity
\[
 \operatorname{Res}_u\operatorname{pp}_u(u^mY(J,u)w)=J_mw
\]
therefore gives
\begin{equation}\label{eq:vol2-current-first}
 H_m=-\operatorname{Res}_u\widetilde h(u^mJ\boxtimes-),
 \qquad d_{\mathrm{Hom}}H_m=J_ma-aJ_m.
\end{equation}
The source action is the rightmost $J_m$ in this formula. In particular,
the zero and negative modes are part of the comparison. Equality of
the restrictions in cohomology would not specify $\gamma,\delta$ or
the corrected supported coordinate.

\subsection{The second equation and all higher equations}

Let $\mathfrak a$ be an ordinary Lie algebra and let $M,V$ be dg
$\mathfrak a$-modules. Put $E=\operatorname{Hom}_k(M,V)$ and
$T_xf=\rho_V(x)f-f\rho_M(x)$. For an alternating $p$-cochain set
\begin{align}\label{eq:vol2-current-ce}
 (\delta f)(x_0,\ldots,x_p)
 ={}&\sum_i(-1)^iT_{x_i}f(x_0,\ldots,\widehat x_i,\ldots,x_p)
 \notag\\
 &+\sum_{i<j}(-1)^{i+j}
 f([x_i,x_j],x_0,\ldots,\widehat x_i,\ldots,\widehat x_j,\ldots,x_p).
\end{align}
The bracket occupies the first slot. The representation identity and
Jacobi identity give $\delta^2=0$; moreover $\delta d_E=d_E\delta$.
Use the product totalization with differential
$\delta+(-1)^p d_E$ on Chevalley degree $p$.

A coherent current-linear map is a family
$H_p:\Lambda^p\mathfrak a\to E^{-p}$ satisfying
\begin{equation}\label{eq:vol2-current-hierarchy}
 H_0=a,\qquad d_Ea=0,\qquad d_EH_p=\delta H_{p-1}\quad(p\geq1).
\end{equation}
Indeed, $\alpha_p=(-1)^{p(p-1)/2}H_p$ is then a total degree-zero
cocycle. At two inputs the equation is
\begin{equation}\label{eq:vol2-current-second}
 d_EH_2(x,y)=T_xH_1(y)-T_yH_1(x)-H_1([x,y]).
\end{equation}
Its right side is closed once the first equation holds. Its class in
$\operatorname{Hom}(\Lambda^2\mathfrak a,H^{-1}E)$ must vanish
before this chosen first homotopy extends. The three terms keep the
two orders of current--root collision and the current--current
collision. The last is the third channel.

\begin{theorem}[A resolved current action]\label{thm:vol2-resolved-current}
Suppose $I$ is a dg $\mathfrak a$-module and $i:V\to I$ is equivariant.
Let $(p,i,h)$ be a contraction of complexes with
\[
 pi=1,\qquad d_Ih+hd_I=1-ip,\qquad hi=ph=h^2=0.
\]
For every equivariant cochain map $j:M\to I$ there is a coherent
current-linear map with unary part $pj$. It is given in all arities
by a series finite in each output arity. Neither $p$ nor $h$ need
be equivariant.
\end{theorem}
\begin{proof}
Postcomposition gives the same contraction on the Hom complexes.
On Chevalley degree $p$ its homotopy is $h_C=(-1)^ph$. Treat
$\delta$ as a perturbation of $(-1)^pd_{\mathrm{Hom}}$ and define
\begin{equation}\label{eq:vol2-current-transfer}
 p'=p(1+\delta h_C)^{-1}
    =\sum_{r\geq0}(-1)^rp(\delta h_C)^r.
\end{equation}
Each perturbation raises Chevalley degree, so this defines a map on
the product totalization. Since $i$ is equivariant and $hi=0$,
$h_C\delta i=0$. The transferred differential is consequently
the ordinary total Chevalley differential on $E$.
Multiplying the geometric series and using the contraction identity
shows that $p'$ intertwines the total differentials. The map $j$ is
a total cocycle in arity zero. Set
$H_p=(-1)^{p(p-1)/2}(p'j)_p$ to obtain
\eqref{eq:vol2-current-hierarchy}.

For the first two coefficients the formulas reduce to
\begin{align}\label{eq:vol2-current-transfer-first}
 H_1(x)&=-p\rho_I(x)hj,\\
 H_2(x,y)&=p\bigl(\rho_I(x)h\rho_I(y)h
                      -\rho_I(y)h\rho_I(x)h\bigr)j.
 \label{eq:vol2-current-transfer-second}
\end{align}
For example, applying $d_E$ to the first gives
$-p\rho_I(x)(1-ip)j=\rho_V(x)pj-pj\rho_M(x)$.
In the second, differentiation of the two copies of $h$ leaves
the two action terms and their commutator, giving
\eqref{eq:vol2-current-second}. The terms containing $ph$, $hi$
or $h^2$ vanish by the side conditions.
\end{proof}

\begin{example}\label{ex:vol2-second-current}
The second homotopy may be nonzero for commuting currents. Let $I$
have basis $z,c$ in degree $-2$, $d,a$ in degree $-1$, and $b$ in
degree zero, with $d_Ic=d$ and $d_Ia=b$. Put $V=kz$, let $p$ be
projection, and take $h(d)=c$, $h(b)=a$. Define commuting cochain
operators by $\rho_x(c)=z$, $\rho_y(a)=d$, with all other values
zero. For $M=I$ and $j=1$ the formulas give
\[
 H_1(x)(d)=-z,\qquad H_1(y)=0,\qquad H_2(x,y)(b)=z.
\]
All other values vanish. Thus $d_EH_2(x,y)(a)=-z$, which is
the nonzero right side of \eqref{eq:vol2-current-second} on $a$.
\end{example}

For an infinite current algebra the same construction requires either
pointwise finiteness of the operations or a specified completed
Chevalley domain. The theorem begins with an equivariant resolved
presentation. Applying it to a geometric boundary action still
requires a map from that action's supported collision complex to
this presentation, preserving its unary path and all three channels.

\section{Several roots and upward extensions}
\label{sec:vol2-root-support}

Return to the Heisenberg state space $V=V_\kappa$ of
\eqref{eq:vol2-fock-currents}. The acting current adds a root to
the source of a module cochain.
This explains why a bound on the roots of the original module is not
a bound on every output of its derived endomorphisms. We make the
coefficient category explicit before proving the bound that does hold.

For a finite set $I$, take a chart $X^I_\sigma$ for each total order
$\sigma$ of $I$, where $X=\mathbb A^1_k$. Glue two charts along the
open set obtained by removing $x_i=x_j$ whenever their orders of
$i,j$ disagree. Denote the result by $\mathscr R_I$. The triple
overlap condition holds: disagreement of two orders on a pair forces
one of them to disagree with any third order on that pair. Thus these
opens define a smooth, generally nonseparated scheme. At a tuple
with equality blocks $B$, its fibre over $X^I$ has
$\prod_B|B|!$ points, one order for each block.

A surjection $\alpha:I\twoheadrightarrow J$ with an order on each
fibre defines $\delta_\alpha:\mathscr R_J\to\mathscr R_I$ by
$x_i=y_{\alpha(i)}$: expand every ordered label of $J$ into its
ordered fibre. Composition concatenates these orders. The maps
glue because the preimage of a chart overlap is the corresponding
overlap on $J$. Let $\mathrm{OSurj}_+$ be the category of nonempty
finite sets and these ordered-fibre surjections. Define
\begin{equation}\label{eq:vol2-full-coefficient-category}
 \mathcal C_+=\varprojlim_{\mathrm{OSurj}_+}
                 D\operatorname{Mod}(\mathscr R_I),
 \qquad \mathcal C=D(k)\times\mathcal C_+,
\end{equation}
using extraordinary inverse images $\delta_\alpha^!$ on right
differential modules. Differential modules on the atlas are defined
by descent from its charts. An object retains the equivalences
$\delta_\alpha^!F_I\simeq F_J$ and all their coherences.
The factor $D(k)$ is the empty coefficient; a positive object has
zero empty coefficient.

For a decomposition $I=A\sqcup B$ let $U_{A,B}$ be the locus on
which every cross pair is distinct, and let $j_{A,B}$ be its open
inclusion. Restriction of orders identifies this open with the
disjoint locus in $\mathscr R_A\times\mathscr R_B$. Set
\begin{equation}\label{eq:vol2-full-convolution}
 (F\star G)_I=
 \bigoplus_{I=A\sqcup B}j_{A,B,*}
             (F_A\boxtimes G_B)|_{U_{A,B}}.
\end{equation}
An empty subset contributes scalar tensor product with its coefficient
in $D(k)$; at $I=\varnothing$ the value is
$F_\varnothing\otimes_kG_\varnothing$.
Both subsets may be empty. This is the all-subset convolution; the
cut product of Section~\ref{sec:vol2-word-coefficients} used only
consecutive decompositions of a fixed word.

\begin{proposition}\label{prop:vol2-full-convolution}
The product \eqref{eq:vol2-full-convolution} is symmetric monoidal.
Its unit $K$ is $k$ at the empty coefficient and zero elsewhere.
Every parenthesized product of $r$ positive objects has the common
formula
\begin{equation}\label{eq:vol2-full-partitions}
 (F_1\star\cdots\star F_r)_I
 =\bigoplus_{I=I_1\sqcup\cdots\sqcup I_r}
      j_{P,*}(F_{1,I_1}\boxtimes\cdots\boxtimes F_{r,I_r})|_{U_P},
\end{equation}
where $P=(I_1,\ldots,I_r)$ and $U_P$ removes all cross-block
diagonals. In particular this coefficient is zero for $r>|I|$.
\end{proposition}
\begin{proof}
Under an ordered collision, a summand whose blocks meet the same
collapsed fibre has zero extraordinary inverse image: that fibre
misses its open locus. Every remaining block is the inverse image
of a unique block downstairs. Open base change
(Lemma~\ref{lem:vol2-open-basechange}) identifies these remaining
summands with the convolution downstairs. This also proves
compatibility with composed collisions.

Nested open conditions impose exactly the cross-block inequalities
of $U_P$. Localization commutes with external products, so all
parenthesizations have \eqref{eq:vol2-full-partitions} as their
common expression. The coherent base-change identifications give
the associativity diagrams. Permutation of the blocks gives the
symmetry with the usual Koszul signs. Only an empty block can
contribute the object $K$, proving the unit assertion. Positivity
forces every $I_j$ to be nonempty, proving the last assertion.
\end{proof}

For a right differential module $M$ on $X$, its main-diagonal family
$iM$ is $\Delta_{I,*}M$ on each ordered chart. The branch maps are
the extraordinary restrictions of this diagonal direct image.
Let $V^r=\omega_X\otimes_O\mathcal V$. The residue pairings of the
Heisenberg fields give the degree-zero Lie object
\[
 \mathfrak g=L_{V_\kappa}=iV^r
\]
for $\star$. Explicitly, on a two-block summand its bracket is the
principal part of $Y(v,z-y)w$, with its right transfer-module
normal derivatives. The field skew symmetry and Jacobi identity
are the antisymmetry and Jacobi identities of this pairing, on
every ordered branch. The local vacuum is the state $1\in V$;
it is distinct from the empty convolution unit $K$.

Write
\[
 N_a=\operatorname{Sym}^a_\star\mathfrak g,\qquad
 A=U_\star\mathfrak g.
\]
Symmetric powers use coinvariants, identified with invariants by
averaging. We use the full derived category of $A$-modules and
the adjoint action on $N_a$ and $A$, without imposing preservation
of conformal weight. PBW symmetrization gives
$A=\bigoplus_{a\geq0}N_a$ as adjoint modules. The usual multilinear
PBW identities apply here: convolution preserves colimits, and
finite symmetric-group coinvariants are exact in characteristic
zero. Since $N_a(I)=0$ for $a>|I|$, this sum is also the product
in compatible coefficient families. On each chart, localization
and diagonal direct image are exact, so each $N_a$ has
differential-module coefficients in degree zero.

\begin{theorem}[The degree-dependent support bound]
\label{thm:vol2-root-bound}
For $w,a\geq1$,
\begin{equation}\label{eq:vol2-root-bound}
 \operatorname{Ext}^q_{\mathfrak g}(N_w,N_a)=0\qquad(q<a-w).
\end{equation}
Consequently, for $w\geq1$ and $q\geq0$,
\begin{equation}\label{eq:vol2-root-window}
 \tau_{\leq q}R\operatorname{Hom}_{\mathfrak g}(N_w,A)
 \simeq\tau_{\leq q}\prod_{0\leq a\leq w+q}
              R\operatorname{Hom}_{\mathfrak g}(N_w,N_a).
\end{equation}
The truncation is applied to the full complexes. In a connective
coefficient resolution, degree-$q$ cycle equations retain Chevalley
columns through $q+1$.
\end{theorem}
\begin{proof}
The coefficient family $N_a$ is a sum of open direct images from
exactly $a$ distinct roots, with diagonal support within each
block. If a coefficient family $B$ is supported on at most
$a-1$ roots, ordinary open adjunction gives
$R\operatorname{Hom}_{\mathcal C}(B,N_a)=0$.
This is an adjunction before imposing module equivariance.

The module Chevalley resolution has columns
\[
 R\operatorname{Hom}_{\mathcal C}
       (\Lambda^n_\star\mathfrak g\star N_w,N_a)[-n].
\]
Their sources have at most $n+w$ roots, so a column vanishes
when $n+w<a$. Otherwise it is $n$-connective: its coefficient
objects are differential modules in degree zero and their
derived Homs have no negative cohomology. The compatible-family
limits preserve this bound. In each total degree only finitely
many columns contribute. This proves \eqref{eq:vol2-root-bound}.

Derived Hom into $A$ is the product of derived Homs into the
$N_a$. Every omitted factor in \eqref{eq:vol2-root-window} starts
above degree $q$; products preserve that lower bound. Finally
column $q+1$ can receive the outgoing differential of a
degree-$q$ cochain, whereas later columns cannot enter its
cycle equation.
\end{proof}

The boundary of this vanishing range is attained. Let
$\mathsf D=\partial_{x_1}$ on polynomial states. It satisfies
$\mathsf D1=0$, $\mathsf DJ=1$ and $[\mathsf D,T]=0$.
Moreover $[\mathsf D,Y(J,u)]=1=Y(\mathsf DJ,u)$.
Normal products and translation extend this derivation identity
to every state. Thus $\phi_\tau=e^{\tau\mathsf D}$ is a
vertex-algebra automorphism over $k[\tau]$, finite on each state,
and induces a Lie derivation of $\mathfrak g$. This is the current
translation subgroup described by Matsuo--Nagatomo
\cite[\S3.3, Theorem 3.6]{MatsuoNagatomo1997}.

\begin{proposition}\label{prop:vol2-upward-current}
Let $m_w:\mathfrak g\star N_w\to N_{w+1}$ be symmetric
multiplication. For every $w\geq1$,
\[
 \gamma_w(x;m)=m_w(\mathsf Dx,m)
\]
defines a nonsplit extension of $N_w$ by $N_{w+1}$ as
$\mathfrak g$-modules, preserving the vacuum action.
\end{proposition}
\begin{proof}
The Chevalley differential of $\gamma_w$ is
\[
 m_w([x,\mathsf Dy]+[\mathsf Dx,y]-\mathsf D[x,y],m)=0.
\]
Hence the action on $N_{w+1}\oplus N_w$ is
$x(n,m)=(xn+\gamma_w(x;m),xm)$. Its off-diagonal vacuum action
vanishes because $\mathsf D1=0$.

Ordinary coefficient adjunction gives
$\operatorname{Hom}_{\mathcal C}(N_w,N_{w+1})=0$, so the underlying
split extension has exactly one coefficient splitting. This
splitting is not equivariant: at $w+1$ distinct roots,
$\gamma_w$ sends $J$ at the acting root and vacua at the other
roots to a nonzero external vacuum. Degree-zero maps between
these degree-zero modules are ordinary equivariant maps by the
nonnegative Chevalley resolution. The extension therefore remains
nonzero in the derived category.
\end{proof}

An unrestricted derived open-support adjunction for these modules
would force this extension to vanish. Such an adjunction cannot be
inferred from the corresponding adjunction for their underlying
coefficient families. For instance, input weight two has a genuine
degree-one output in $N_3$.

\subsection{The extension integrates to an interface}

For any degree-zero Lie object with a natural derivation $\mathsf D$,
the formula
\begin{equation}\label{eq:vol2-triangular-current-action}
 \rho_{0,\tau}(x)m=\operatorname{ad}_x(m)+\tau(\mathsf Dx)m
\end{equation}
defines a module structure on $\operatorname{Sym}_\star\mathfrak g[\tau]$.
Indeed, if $M_z$ is multiplication by $z$, then
\[
 [\operatorname{ad}_x,M_{\mathsf Dy}]=M_{[x,\mathsf Dy]},\qquad
 [M_{\mathsf Dx},\operatorname{ad}_y]=M_{[\mathsf Dx,y]},\qquad
 [M_{\mathsf Dx},M_{\mathsf Dy}]=0.
\]
Their sum is the representation identity. The added term raises
symmetric degree by one. For $w\geq1$ and $r\geq0$, quotienting
the degrees at least $w$ by
those above $w+r$ therefore gives a flat family of modules on
$\bigoplus_{a=w}^{w+r}N_a$. For $r=1$ its extension coordinate is
$\tau\gamma_w$.

Suppose $\mathsf D$ is locally nilpotent and the Rees envelope
$A_\nu=U_{\nu,\star}\mathfrak g$, with
$xy-yx=\nu[x,y]$, is PBW flat. Write $\phi_t=e^{t\mathsf D}$.
PBW gives a unique integral operator $\operatorname{ad}_\nu(x)$
with $\nu\operatorname{ad}_\nu(x)m=xm-mx$. Then
\begin{equation}\label{eq:vol2-rees-current-action}
 \rho_{\nu,\tau}(x)m
 =\operatorname{ad}_\nu(x)m+
   \sum_{j\geq1}\frac{\nu^{j-1}\tau^j}{j!}(\mathsf D^jx)m
\end{equation}
is integral over $k[\nu,\tau]$. After inverting $\nu$ it equals
$\nu^{-1}(\phi_{\nu\tau}(x)m-mx)$.
Commuting left and right multiplications gives
\[
 [L_{\phi x}-R_x,L_{\phi y}-R_y]
   =\nu(L_{\phi[x,y]}-R_{[x,y]}).
\]
Torsion-freeness then proves the representation identity before
inversion. At $\nu=0$ only the $j=1$ shifted term survives, giving
\eqref{eq:vol2-triangular-current-action}.

For an associative algebra $A$ and its automorphisms $\phi_\tau$,
let $B_\tau={}_{\phi_\tau}A$ have ordinary right action and left
action $x\cdot b=\phi_\tau(x)b$. These twisted regular bimodules
compose by
\begin{equation}\label{eq:vol2-interface-composition}
 B_\sigma\otimes_A^L B_\tau\xrightarrow{\sim}B_{\sigma+\tau},
 \qquad a\otimes b\longmapsto\phi_\tau(a)b.
\end{equation}
They are free on either side, so ordinary tensor computes the
derived tensor. The formula is balanced and left equivariant;
its inverse is $c\mapsto1\otimes c$. Both associations of three
factors give $\phi_{\tau+\upsilon}(a)\phi_\upsilon(b)c$.
Thus $B_0=A$ is the unit and $B_{-\tau}$ is inverse to $B_\tau$.
Equation~\eqref{eq:vol2-rees-current-action} is the normalized
left-minus-right action on $B_{\nu\tau}$. This constructs the
algebraic interface underlying the upward extension. A field
realization must also supply its collision and interchange maps.

Choose semi-free resolutions $P_w\to N_w$ and form the dg category
$\mathcal E(w,a)=\operatorname{Hom}_{\mathfrak g}(P_w,P_a)$.
Composition is composition of resolution maps. The dg Yoneda
functor identifies the subcategory generated by the $P_w$ under
finite cones, shifts and retracts with $\operatorname{Perf}(\mathcal E)$:
it is fully faithful on the generators, this property extends by
those operations in each variable, and represented modules generate
the perfect modules. This is the precise generated-sector conclusion
of dg Morita theory \cite[Theorem 3.8(b)]{KellerDG2006}.
It supplies no generation statement for all physical boundaries.

\section{Operators at a polynomial root}
\label{sec:vol2-polynomial-operators}

Forgetting the current action leaves an endomorphism of the moving
differential module. The allowed endomorphisms must be specified
before this forgetful map can be computed. Put
\begin{equation}\label{eq:vol2-affine-operators}
 N=\operatorname{End}_O(\mathcal V),\qquad
 \nabla_NA=\nabla A-A\nabla,\qquad
 G(N)=[N\xrightarrow{\nabla_N}N],
\end{equation}
where $G(N)$ is a complex of $k$-vector spaces in degrees zero and
one. The identity $\nabla_N(fA)=f'A+f\nabla_NA$ for $f\in O$
makes $N$ a differential module. In a monomial
state basis, $N$ consists of column-finite matrices over $O$: each
input has a finite polynomial output, with no common bound on all
outputs. For example, $v_j\mapsto y^jv_j$ lies in $N$ but not in
$O\otimes_k\operatorname{End}_k(V)$.

\begin{lemma}\label{lem:vol2-internal-operator}
For an $O$-free left $D_y$-module $M$ in degree zero, the internal
derived endomorphism for $\otimes_O^L$ is
$\operatorname{End}_O(M)$ in degree zero, with its commutator
connection. This holds without a finite-rank hypothesis.
\end{lemma}
\begin{proof}
Evaluation gives
\[
 \operatorname{Hom}_{D_y}(P,\operatorname{Hom}_O(M,Q))
 =\operatorname{Hom}_{D_y}(P\otimes_OM,Q).
\]
The connection on the Hom is forced by this equality. Tensoring
with $M$ is exact. Its right adjoint is exact as well: after choosing
an $O$-basis, its underlying module is a product of copies of $Q$,
and products of modules are exact. Both functors preserve
quasi-isomorphisms, so their adjunction descends to derived
categories. Taking $Q=M$ proves the assertion.
\end{proof}

This affine internal Hom need not commute with localization of an
infinite-rank source. The supported calculation below uses the
restriction of a specified coefficient object, not a new internal
Hom formed on the open set.

\subsection{The supported degrees}

Use right densities $dy$ on $X$ and $dy\wedge du$ on $Y=X^2$.
The monoidal left-to-right normalization for the diagonal tensor is
$M\mapsto M^r[1]$, with $M^r=\omega_X\otimes_OM$; its units are
$J_X=\omega_X[1]$ and $J_Y=\omega_Y[2]$.
Subscripts $\mathcal D$ on Hom denote right differential modules.
For any left $D_y$-module $L$ with connection $\nabla_L$, write
$G(L)=[L\xrightarrow{\nabla_L}L]$ in degrees zero and one.
Its normal direct image has the model
\[
 \Delta_+L=L[\eta],\qquad \partial_u=\eta,\qquad
 u=-\partial_\eta.
\]
Its normal Spencer complex is
$[L[\eta]\xrightarrow{\eta}L[\eta]]$ in degrees zero and one.
Project the degree-one term to its constant coefficient and include
that coefficient back. The homotopy sends $\eta^m\ell$ in degree
one to $\eta^{m-1}\ell$ for $m\geq1$ and to zero for $m=0$.
It satisfies $dh+hd=1-ip$ on every monomial. In tangential form
degree $a$, give both the normal differential and this homotopy
the sign $(-1)^a$. The tangential terms then cancel, proving
\begin{equation}\label{eq:vol2-affine-supported-hom}
 R\operatorname{Hom}_{\mathcal D_Y}
       (J_Y,\Delta_*(L^r[1]))\simeq G(L)[-2].
\end{equation}
One shift comes from the normal Spencer complex, and the other is
the density shift $[1-2]$.

In particular, suppose a coefficient object $E_2$ has
$\Delta^!E_2=N^r[1]$. Recollement and open adjunction give
\begin{equation}\label{eq:vol2-affine-restriction-fibre}
 \operatorname{Fib}\left(
 R\operatorname{Hom}_{\mathcal D_Y}(J_Y,E_2)
 \longrightarrow R\operatorname{Hom}_{\mathcal D_U}(J_U,j^!E_2)
 \right)\simeq G(N)[-2].
\end{equation}
The fibre has cohomology only in degrees two and three. Restriction
therefore induces an equivalence on the unshifted mapping spaces.
For a source $J_Y[-d]$, the fibre instead is $G(N)[d-2]$.
Thus shifted sources detect supported operators invisible to the
unshifted comparison.

For example, let $P_0$ be projection to the constant state.
Both $P_0$ and $\mathsf D=\partial_{x_1}$ commute with $T$ and are
horizontal endomorphisms. Their commutator is nonzero, since
$(P_0\mathsf D-\mathsf DP_0)J=1$. They give classes in the
degree-zero supported fibre for $d=2$. The inclusion of that fibre
need not be injective on cohomology; realizing the operators by
full boundary actions is an additional question.

\subsection{Translation towers}

Give $x_n$ energy $n$ and let $V_+$ be the positive-energy states.
If $x_N$ is the largest variable occurring in a polynomial $f$,
the coefficient of $x_{N+1}$ in $Tf$ is
$N\partial_{x_N}f$. Descent on $N$ proves $\ker T=k1$.
Choose a graded complement $U$ of $TV_+$ in $V_+$, containing
$J$ and $J^2$. Induction on energy gives
\begin{equation}\label{eq:vol2-root-towers}
 V=k1\oplus k[T]\otimes_kU,\qquad
 \mathcal V\cong O\oplus P,\qquad P=D_y\otimes_kU.
\end{equation}
The second map sends $f(y)\partial_y^n\otimes u$ to $f(y)T^nu$;
PBW and the connection identify the two sides. If $p(E)$ is the
partition number, $\dim U_1=1$ and
$\dim U_E=p(E)-p(E-1)$ for $E\geq2$.

On one tower write $t$ for $T$, so its connection is
$\partial_y+t$ on $k[y,t]$. Define
\begin{align}\label{eq:vol2-tower-reduction}
 \operatorname{red}\left(\sum_n f_n(y)t^n\right)
    &=\sum_n(-1)^nf_n^{(n)}(y),\\
 H_P(f(y)t^n)&=\sum_{j=0}^{n-1}(-1)^jf^{(j)}(y)t^{n-1-j},
 \qquad H_P(f)=0.\label{eq:vol2-tower-homotopy}
\end{align}
If $j$ includes the constant tower coefficient, cancellation gives
\[
 \operatorname{red}\nabla=0,\qquad
 \nabla H_P=1-j\operatorname{red},\qquad H_P\nabla=1.
\]
On the vacuum use evaluation at $y=0$ and
$H_O(f)=\int_0^yf(s)\,ds$. Together these maps define a contraction
\begin{equation}\label{eq:vol2-state-contraction}
 G(\mathcal V)\rightleftarrows k\oplus W[-1],
 \qquad W=k[y]\otimes_kU.
\end{equation}
Write $H_{\mathcal V}$ for its homotopy and
$\operatorname{red}_+$ for reduction of its positive-state summand.

Evaluation at the vacuum splits $N$ as differential modules:
\[
 N=\sigma(\mathcal V)\oplus N_0,
 \qquad \sigma(v)=vP_0,\qquad N_0=\{A:A1=0\}.
\]
For $A\in N_0$ define an $O$-linear operator $R(A)$ by
\begin{equation}\label{eq:vol2-column-recursion}
 R(A)1=R(A)u=0,\qquad
 R(A)(T^{n+1}u)=\nabla R(A)(T^nu)-A(T^nu).
\end{equation}
Every column requires finitely many operations, so $R(A)$ is an
allowed endomorphism. If $F\in N_0$, let $\Phi(F)$ be the
horizontal operator with $\Phi(F)(T^nu)=\nabla^nF(u)$ and
$\Phi(F)1=0$. Recursion then gives
\[
 \nabla_NR(A)=A,\qquad R(\nabla_NF)=F-\Phi(F).
\]

\begin{theorem}\label{thm:vol2-polynomial-operators}
The operator complex has a contraction
\[
 G(N)\rightleftarrows B_0\oplus W[-1],
 \qquad B_0=\operatorname{End}_{D_y}(\mathcal V).
\]
Its degree-one projection is $A\mapsto\operatorname{red}_+(A1)$,
and its homotopy is
\begin{equation}\label{eq:vol2-operator-homotopy}
 H_N(A)=\sigma(H_{\mathcal V}(A1))+R(A-\sigma(A1)).
\end{equation}
For a basis $\mathcal U$ of $U$, in the order $(O,P)$ one has
\begin{equation}\label{eq:vol2-horizontal-operator-ring}
 B_0=\begin{pmatrix}
 k&\displaystyle\prod_{u\in\mathcal U}k[y]\\
 0&\operatorname{Mat}_{\mathrm{col},\mathcal U}(D_y^{\mathrm{op}})
 \end{pmatrix}.
\end{equation}
The complex computes $R\operatorname{Hom}_{D_y}(\mathcal V,\mathcal V)$.
\end{theorem}
\begin{proof}
The preceding contractions apply to the two summands of $N$.
The degree-zero projection is the horizontal operator whose value
on the vacuum is the constant-state coefficient of $A1$ at $y=0$
and whose value on each primitive $u$ is $A(u)$. The degree-one
section is $w\mapsto\sigma(jw)$. The displayed identities prove
$dH_N+H_Nd=1-ip$, $pi=1$, $H_Ni=pH_N=H_N^2=0$.

The free module $P$ has arbitrary images on its basis $U$.
The connection $\partial_y+t$ is injective on a polynomial tower,
so $\operatorname{Hom}_{D_y}(O,P)=0$. A map $O\to O$ is a
scalar; a map $P\to O$ chooses a polynomial for each generator.
Endomorphisms of the direct sum of copies of $D_y$ are the
column-finite matrices over $D_y^{\mathrm{op}}$. This proves
\eqref{eq:vol2-horizontal-operator-ring}.

Finally $P$ is projective and
$[D_y\xrightarrow{(-)\partial_y}D_y]$ resolves $O$ in degrees
$-1,0$. Equivalently, the one-dimensional relative Spencer
resolution for the $O$-free module $\mathcal V$ gives the
commutator complex \eqref{eq:vol2-affine-operators}. Thus it is
a model of the stated derived Hom.
\end{proof}

Its cohomology product has $W^2=0$. A matrix
$b=\left(\begin{smallmatrix}a&v\\0&C\end{smallmatrix}\right)$
acts on the left of $W$ through $C$ and on the right by $a$.
On $D_y/\partial_yD_y=k[y]$, right multiplication by
$\partial_y$ acts as $f\mapsto-f'$. This identifies the
cohomology algebra. The contraction was a contraction of complexes;
it has made no assertion about higher compositions.

\subsection{The charge line and the comparison homotopy}

Define the polynomial current cochain algebra
\[
 Q_\kappa=k[y,x_1,x_2,\ldots]\otimes
                  \Lambda(\xi_0,\xi_1,\ldots),\qquad
 dx_n=\kappa n\xi_n,\qquad d\xi_n=0.
\]
The ghosts $\xi_n$ have degree one. Translation extends by
$T\xi_n=(n+1)\xi_{n+1}$. Take its root de Rham complex
$\mathcal C_\kappa=Q_\kappa\oplus dy\,Q_\kappa$, with total
differential $d+dy(\partial_y+T)$ and forms written on the left;
in particular, $d(dy\,q)=-dy\,dq$. Ghost deletion and multiplication
define the cochain map
\begin{equation}\label{eq:vol2-current-forgetting}
 \rho_\kappa:\mathcal C_\kappa\longrightarrow G(N),
 \qquad f\longmapsto L_f,
\end{equation}
with every positive ghost monomial sent to zero. This is the
coefficient map that forgets the regular-current equivariance.
It keeps the differential-module connection.

\begin{proposition}\label{prop:vol2-charge-image}
The current complex has cohomology $k$ in degree zero and $k\gamma$
in degree one, where
\[
 \gamma=\xi_0+dy\,J/\kappa.
\]
The maps induced by $\rho_\kappa$ are
$k\hookrightarrow B_0$, $a\mapsto a\operatorname{id}$, and
$k\gamma\hookrightarrow W$, $\gamma\mapsto J/\kappa$.
\end{proposition}
\begin{proof}
Contract every Koszul pair $(x_n,\xi_n)$, $n>0$, by the Euler
homotopy: on positive total polynomial-plus-ghost degree divide
$\sum_{n>0}x_n(\kappa n)^{-1}\partial_{\xi_n}$ by that degree.
The remaining complex is $k[y]\otimes\Lambda(\xi_0)$.
Perturb by root translation. The perturbation series terminates
after one root form, since $dy^2=0$, and the induced translation
on the zero ghost vanishes because $T\xi_0=\xi_1$.
The transferred differential is $dy\partial_y$, while the
corrected ghost representative is $\gamma$. Directly, its two
differentials are $dy\xi_1$ and $-dy\xi_1$.
Polynomial integration in $y$ leaves the two stated classes.
Finally $\rho_\kappa\gamma=dy\,L_J/\kappa$, whose normal form
is $J/\kappa$ by Theorem~\ref{thm:vol2-polynomial-operators}.
\end{proof}

\begin{corollary}\label{cor:vol2-charge-lift}
For $A\in N$, the class of $dy\,A$ lifts through
$\rho_\kappa$ exactly when $\operatorname{red}_+(A1)\in kJ$.
If this normal form is $cJ/\kappa$, a lift is $c\gamma$, with
comparison homotopy $H_N(A-cL_J/\kappa)$. The fibre satisfies
\[
 H^1\operatorname{Fib}(\rho_\kappa)=B_0/k\operatorname{id},
 \qquad H^2\operatorname{Fib}(\rho_\kappa)=W/kJ,
\]
and has no other cohomology.
\end{corollary}
\begin{proof}
The image criterion follows from Proposition~\ref{prop:vol2-charge-image}.
When it holds, the degree-one projection of $A-cL_J/\kappa$
vanishes, so the contraction identity gives
$\nabla_NH_N(A-cL_J/\kappa)=A-cL_J/\kappa$.
The two injections in that proposition give the fibre groups.
\end{proof}

For example, $L_{J^2}$ and $yL_J$ do not lift: their normal forms
are $J^2$ and $yJ$. A scalar charge $q$ does integrate, with
$J_0=q\operatorname{id}$ and
$\nabla_q=\partial_y+T+(q/\kappa)L_J$. Only the $n=1$ moving-mode
equation changes, by
$[(q/\kappa)L_J,\kappa\partial_{x_1}]=-q$.
When $q\ne0$ the vector $1$ is no longer horizontal.

There are two distinct completions. For a new parameter $g$,
suppose $A(g)\in gN[[g]]$ kills $1$. The map
\[
 U_g1=1,\qquad U_g(T^nu)=(\partial_y+T+A(g))^nu
\]
extends $k[y][[g]]$-linearly and intertwines the two connections.
It is the identity modulo $g$, so it has a $g$-adic inverse.
Every coefficient of every column is a polynomial. This trivializes
the connection with a chosen horizontal vacuum, but leaves its
horizontal automorphisms as stabilizers.

Completing the root instead gives $\widehat{\mathcal V}=V[[y]]$.
The gauge $e^{-yT}$ turns its connection into $\partial_y$; on
continuous endomorphisms it acts by conjugation. Formal integration
is onto, hence
\begin{equation}\label{eq:vol2-completed-root-operators}
 H^0G(\widehat N)=\operatorname{End}_k(V),\qquad
 H^1G(\widehat N)=0,
 \qquad
 \widehat N=\operatorname{End}_{k[[y]],\mathrm{cont}}(V[[y]]).
\end{equation}
In this category the charge has the primitive
\[
 g_\kappa(y)=\frac1\kappa\sum_{n\geq0}
          \frac{(-1)^ny^{n+1}}{(n+1)!}T^nJ,
 \qquad (\partial_y+T)g_\kappa=J/\kappa.
\]
Its infinitely many nonzero state coefficients prevent it from
being a polynomial-root primitive. Completion has changed the
coefficient category and the cohomology.

\section{A nonformal derived operator algebra}
\label{sec:vol2-root-composition}

One free translation tower and the vacuum already detect the
composition missing from the linear contraction. Write $D=D_y$
and choose the tower generated by $J$. Let $P=D$ and
$O=D/D\partial$. Resolve $O$ by
\[
 Q=[D\xrightarrow{R_\partial}D],\qquad
 Q^{-1}=D,\quad Q^0=D,\quad R_\partial(a)=a\partial.
\]
The Weyl PBW basis proves injectivity of $R_\partial$ and identifies
its cokernel with $O$. Put
\[
 B=D^{\mathrm{op}}=k\langle y,p\rangle/(py-yp+1),\qquad
 \mathscr E=\operatorname{End}_D(P\oplus Q).
\]
Thus $p$ represents right multiplication by $\partial$, with
$[p,y]=-1$. Products in $\mathscr E$ are matrix compositions, the
rightmost map acting first. The free cells $P,Q^0,Q^{-1}$ have
degrees $(0,0,-1)$; the only nonzero entry of their differential
is $(Q^0,Q^{-1})=p$. Hence
\[
 \delta X=dX-(-1)^{|X|}Xd
\]
makes $\mathscr E$ a dg algebra of ranks $2,5,2$ over $B$ in
degrees $-1,0,1$. It models the marked derived endomorphisms of
$P\oplus O$. The splitting \eqref{eq:vol2-root-towers} embeds it
as a corner in a projective model of the full Heisenberg module.

\subsection{Polynomial division gives a contraction}

Use normal monomials $y^ap^b$. Define
\begin{align}\label{eq:vol2-weyl-remainders}
 r_L(y^ap^b)&=\partial_y^b(y^a),&
 \ell(y^ap^b)&=\sum_{j=0}^{\min(a,b-1)}
                    \frac{a!}{(a-j)!}y^{a-j}p^{b-1-j},\\
 r_R\left(\sum_bf_b(y)p^b\right)&=f_0(y),&
 r\left(\sum_bf_b(y)p^b\right)&=\sum_{b\geq1}f_b(y)p^{b-1}.
 \notag
\end{align}
Empty sums are zero. The relation $yp=py+1$ gives
\begin{equation}\label{eq:vol2-weyl-divisions}
 F=p\ell(F)+r_L(F)=r(F)p+r_R(F).
\end{equation}
In particular $pB=\ker r_L$ and $Bp=\ker r_R$, and neither
contains $1$.

The block $\operatorname{Hom}(P,Q)$ is
$[B\xrightarrow{p(-)}B]$ in degrees $-1,0$. Its projection is
$r_L$, its section includes polynomials, and its homotopy is
$\ell$. The block $\operatorname{Hom}(Q,P)$ is
$[B\xrightarrow{-(-)p}B]$ in degrees $0,1$; its maps are
$r_R$, polynomial inclusion and $-r$.

On $\operatorname{End}(Q)$ write a degree-zero map as $(a,c)$,
acting on $Q^0,Q^{-1}$, and a degree-one map as $z:Q^{-1}\to Q^0$.
Then
\begin{equation}\label{eq:vol2-vacuum-block}
 \delta h=(ph,hp),\qquad \delta(a,c)=pc-ap.
\end{equation}
With $I(f)=\int_0^yf(t)dt$, its projection to $k$ is
$r_L(a)(0)$, the section is $s\mapsto(s,s)$, and its homotopy is
\begin{equation}\label{eq:vol2-vacuum-block-homotopy}
 H(a,c)=\ell(a),\qquad
 H(z)=(-I(r_Lz),\ell(z)-I(r_Lz)).
\end{equation}
It is zero on degree $-1$.

\begin{proposition}\label{prop:vol2-marked-contraction}
These block maps, and the identity on $\operatorname{End}(P)=B$,
form a contraction of $\mathscr E$ onto its cohomology with
$\pi i=1$, $H^2=Hi=\pi H=0$. In the order $(P,O)$ the result is
\[
 H^0\mathscr E=\begin{pmatrix}B&0\\B/pB&k\end{pmatrix},
 \qquad
 H^1\mathscr E=\begin{pmatrix}0&B/Bp\\0&0\end{pmatrix},
 \qquad H^{\ne0,1}\mathscr E=0.
\]
Every map in this contraction is polynomial on each input.
\end{proposition}
\begin{proof}
The off-diagonal identities are the two divisions in
\eqref{eq:vol2-weyl-divisions}. For the vacuum block, the maps
$(a,c)\mapsto r_L(a)$ and $z\mapsto r_L(z)$ map the complex to
$[k[y]\xrightarrow{-\partial_y}k[y]]$.
A section sends $f$ to $(f,f)$ and $g$ to $g$, since
$pf-fp=-f'$. In the kernel, write $a=pa_1$, $z=pz_1$.
The differential becomes
\[
 B\longrightarrow B\oplus B\longrightarrow B,
 \qquad h\longmapsto(h,hp),\quad(a_1,c)\longmapsto c-a_1p.
\]
The homotopy sends $(a_1,c)$ to $a_1$ and $z_1$ to $(0,z_1)$.
Contract the remaining derivative by $-I$. Composition gives
\eqref{eq:vol2-vacuum-block-homotopy} and
$\delta H+H\delta=1-i\pi$. The identity
$\ell(ap)=\ell(a)p+r_L(a)$ verifies the middle degree directly.
Integration vanishes at zero and $\ell$ kills polynomials, giving
the side conditions. All maps preserve the two object idempotents.
\end{proof}

\subsection{The full Massey coset}

Consider the closed maps
\[
 a:P\to Q,\quad a^0=1;\qquad
 b:Q\to P[1],\quad b^{-1}=1;\qquad
 c:P\to P,\quad c=p.
\]
Their degrees are $0,1,0$, and $ba=0$, while
$cb=\delta h$ for $h:Q^0\to P$, $h=-1$.
For primitives $\delta h=cb$, $\delta k=ba$, use the defining
Massey representative $ha-ck$. Taking $k=0$ gives $-1_P$.

\begin{theorem}\label{thm:vol2-root-massey}
The exact composable Massey coset is
\begin{equation}\label{eq:vol2-root-massey}
 \langle c,b,a\rangle=-1+pB\ \subset H^0\operatorname{End}(P)=B.
\end{equation}
It excludes zero. Both $\mathscr E$ and the full polynomial-root
derived endomorphism algebra of $\mathcal V$ are nonformal.
\end{theorem}
\begin{proof}
A second primitive of $cb$ differs by a degree-zero cocycle in
$\operatorname{Hom}(Q,P)$. There are no such nonzero cocycles,
since its differential is right multiplication by $-p$.
A primitive of $ba=0$ can differ by any element of $B$.
The indeterminacy is therefore exactly $pB$. The remainder map
$r_L$ annihilates this ideal and sends $1$ to $1$, so zero is
excluded.

In a dg algebra with zero differential, a defined triple Massey
product contains zero, by choosing zero primitives. A dg
quasi-isomorphism preserves the Massey coset: it sends defining
systems to defining systems, and surjectivity on cohomology
recovers all changes of primitives, whose images are the two
indeterminacy subspaces. Applying this in both directions along
a quasi-isomorphism zigzag proves nonformality; equivalently one
may use the minimal-model criterion of
\cite[\S\S3.3--3.4]{KellerAInfinity2001}.

For the full module, sandwich a defining system between the
idempotents of the chosen $P\oplus O$ summand. Its $P\to P$
projection is a defining system above. Additional matrix blocks
do not change the projected indeterminacy: the degree-zero
$O\to P$ cohomology is still zero, and the $P\to P$ corner
is $B$. A zero total Massey class would thus give zero in
\eqref{eq:vol2-root-massey}, which is impossible.
\end{proof}

Replacing $a$ by the representative $a_f=f(y)$ gives
$\langle c,b,a_f\rangle=-f(y)+pB$, nonzero for every nonzero
polynomial $f$. The obstruction concerns the whole coset, not a
particular transferred ternary coefficient. At level $\kappa$,
$a$ is the restriction of $J_1/\kappa=\partial_{x_1}$ to the
$J$-tower, $b$ is its charge class $[J]$, and $c$ is its translation.
Thus this composition uses three of the operators already present
in the current comparison.

\subsection{Higher compositions and the ordinary centre}

The polynomial contraction gives all higher operations explicitly.
Let $\mathcal H=H^*\mathscr E$, and let
$s:\mathscr E\to\mathscr E[1]$ have degree $-1$. On the tensor
coalgebra of suspended matrices use
\[
 b_1(sx)=-s\delta x,\qquad
 b_2(sx,sy)=(-1)^{|x|}s(xy),\qquad h_s=-sHs^{-1}.
\]
Extend $b_1,b_2$ as coderivations, and write $i_s,\pi_s$ for
the suspended inclusion and projection. The tensor contraction is
\[
 \mathbf H|_{\mathscr E[1]^{\otimes n}}
 =\sum_{j=1}^n(i_s\pi_s)^{\otimes(j-1)}\otimes h_s
                                \otimes1^{\otimes(n-j)},
\]
with graded tensor signs. Let $i_T,\pi_T$ be the tensor extensions
of $i_s,\pi_s$.
Homological perturbation gives
\begin{align}\label{eq:vol2-root-ainfinity}
 b_{\mathcal H}&=\pi_Tb_2(1+\mathbf Hb_2)^{-1}i_T,\notag\\
 \mathbf I&=(1+\mathbf Hb_2)^{-1}i_T,&
 \mathbf P&=\pi_T(1+b_2\mathbf H)^{-1},&
 \mathbf H'&=\mathbf H(1+b_2\mathbf H)^{-1}.
\end{align}
Each binary merger lowers length, so at input length $n$ these
inverse series have at most $n-1$ merger stages. Expanding the
contraction identity yields $b_{\mathcal H}^2=0$,
$(b_1+b_2)\mathbf I=\mathbf I b_{\mathcal H}$ and
$\mathbf P\mathbf I=1$, with comparison homotopy $\mathbf H'$.
Deconcatenation and the tensor homotopy give the coalgebra
identities. Tensoring over the two scalar idempotents retains
composable strings; the side conditions give strict units.
This constructs a minimal $A_\infty$ category in every arity.
Its deciding value is
\[
 (b_{\mathcal H})_3(sc,sb,sa)=-s1_P:
\]
the merger $ba$ vanishes, while the other merger has homotopy
$h=-1$. The Massey coset prevents removal of all higher
operations by a different minimal model.

Its ordinary bimodule centre has a different answer. The free
module $P=D$ generates $\operatorname{Perf}(D)$ and $Q$ is a
two-term cone of free modules. Thus $P\oplus Q$ is another
generator. The functors
\[
 R\operatorname{Hom}_D(P\oplus Q,-),\qquad
 -\otimes_{\mathscr E}^L(P\oplus Q)
\]
are inverse on the generated perfect categories: evaluation
is an equivalence on the generator and extends by cones,
shifts and retracts. Derived Morita invariance of Hochschild
cochains \cite[\S5.4]{KellerDG2006} gives
\begin{equation}\label{eq:vol2-marked-root-centre}
 \operatorname{HH}^*(\mathscr E)
   \cong\operatorname{HH}^*(D^{\mathrm{op}})
   \cong\operatorname{HH}^*(D)=k.
\end{equation}
For the last equality, resolve the Weyl algebra as a bimodule
by its two-generator PBW Koszul complex. Its cochains have
differential given by $[\partial,-]=\partial_y$ and
$[y,-]=-\partial_\partial$, hence are the polynomial de Rham
complex after a change of one exterior generator. Exactness of
the bimodule resolution follows from its associated graded
polynomial Koszul resolution and the bounded-below PBW filtration.
Euler contraction, divided by positive polynomial-plus-form
degree, leaves only the constant class. This proves
\eqref{eq:vol2-marked-root-centre} without a root completion.

Nonformality of a selected derived operator algebra is compatible
with a scalar ordinary centre of its perfect module sector.
The centre of the full geometric boundary theory requires its
own category and its own action comparison.

\section{The level ring and the forgetful fibre}
\label{sec:vol2-integral-current-fibre}

The nonzero-level calculation divided by $\kappa$. To retain the
singular level, work over $R=k[\lambda]$ and use a separate symbol
for the Koszul differential:
\[
 \mathcal Q=k[x_1,x_2,\ldots]\otimes\Lambda(\xi_0,\xi_1,\ldots),
 \qquad d_0x_n=n\xi_n,\quad d_0\xi_n=0.
\]
Translation is $Tx_n=nx_{n+1}$ and
$T\xi_n=(n+1)\xi_{n+1}$. It commutes with $d_0$ on generators
and hence everywhere. With forms on the left define
\begin{equation}\label{eq:vol2-integral-current-complex}
 \begin{split}
 \mathcal C_\lambda&=\mathcal Q_R[y]\oplus dy\,\mathcal Q_R[y],\qquad
        \mathcal Q_R=R\otimes_k\mathcal Q,\\
 \mathsf D_\lambda(q+dy\,r)
   &=\lambda d_0q+dy\bigl((\partial_y+T)q-\lambda d_0r\bigr).
 \end{split}
\end{equation}
This is the polynomial-root current complex at the universal level.
There are only two root-form degrees, so its totalization needs
no additional convention. The parameter $\lambda$ has degree zero.

\subsection{Translation coinvariants before specialization}

Let $\mathcal Q_+$ be the augmentation ideal. Translation is
injective on each $\mathcal Q_+^q$. To see this, choose the largest
occupied index of an even or an odd variable. Its translate
introduces a new variable of that type, with coefficient its
corresponding partial derivative. No other term can cancel that
new variable. Descending on the occupied indices removes every
variable, including $\xi_0$; only the unit in the full algebra
can be translation-invariant.

Put $K^q=k[y]\otimes\mathcal Q_+^q$. The operator
$\nabla=\partial_y+T$ is also injective on $K^q$: the highest
power of $y$ in a nonzero kernel element would have a nonzero
coefficient killed by $T$. Define the vector spaces
\begin{equation}\label{eq:vol2-integral-coinvariants}
 \mathsf W^q=K^q/\nabla K^q,\qquad
 \overline d[q]=[d_0q],\qquad
 \mathsf B^q=\operatorname{im}(\overline d:
                   \mathsf W^{q-1}\to\mathsf W^q)\quad(q\geq1).
\end{equation}
They are $k$-vector spaces. The quotient is not an $O$-module
quotient, since $[\nabla,y]=1$.

For computation, choose translation primitives in each state energy
and exterior degree. These pieces are finite dimensional, including
the energy-zero ghost. The same tower division as
\eqref{eq:vol2-tower-reduction} gives the finite normal form
\[
 [f(y)T^nu]=(-1)^nf^{(n)}(y)[u].
\]
Apply $d_0$, decompose its result into towers, and reduce to compute
$\overline d$. In degree zero $\mathsf W^0=W$ of
\eqref{eq:vol2-state-contraction}, as vector spaces.

\begin{lemma}\label{lem:vol2-coinvariant-cohomology}
The cohomology of $(\mathsf W,\overline d)$ is the
one-dimensional span $k\cdot[J]$ in degree zero and zero in
positive degrees.
\end{lemma}
\begin{proof}
The Euler homotopy for the positive Koszul pairs is
$\sum_{n>0}(x_n/n)\partial_{\xi_n}$ divided by positive total
degree. It leaves only $\xi_0$ in the augmentation ideal, so
$H^*(K,d_0)=k[y]\xi_0$ in degree one. Translation acts trivially
there, because $T\xi_0=\xi_1=d_0J$. Thus $\nabla$ acts on this
cohomology as $\partial_y$, which is onto with kernel $k\xi_0$.
The exact sequence of complexes
\[
 0\longrightarrow K\xrightarrow{\nabla}K
       \longrightarrow\mathsf W\longrightarrow0
\]
now proves the assertion. The connecting preimage of $\xi_0$ is
$[J]$, since $d_0J=\nabla\xi_0$.
\end{proof}

\begin{theorem}\label{thm:vol2-integral-current}
There is an integral quasi-isomorphism of complexes
\begin{equation}\label{eq:vol2-integral-current-reduction}
 \mathcal C_\lambda\simeq
 R\oplus(\mathsf W\otimes_kR,\lambda\overline d)[-1].
\end{equation}
Consequently
\begin{align}\label{eq:vol2-integral-current-cohomology}
 H^0\mathcal C_\lambda&=R,&
 H^1\mathcal C_\lambda&=R\Gamma_\lambda,
       &\Gamma_\lambda&=\lambda\xi_0+dy\,J,\notag\\
 H^{q+1}\mathcal C_\lambda&=\mathsf B^q\otimes_kR/(\lambda)
       &&(q\geq1).
\end{align}
\end{theorem}
\begin{proof}
The constant-state root complex contracts to $R$. On the positive
part the map $q+dy\,r\mapsto[r]$ has target differential
$-\lambda\overline d$, the sign of the shift. In its kernel
$r=\nabla s$ has a unique polynomial solution $s$. The homotopy
sends $dy\nabla s$ to $s$ and the root-degree-zero part to zero.
The terms $\lambda d_0s$ cancel in $\mathsf D_\lambda H+
H\mathsf D_\lambda$, leaving the identity. This proves the
quasi-isomorphism without inverting the level.

Multiplication by $\lambda$ is injective on each extended vector
space. Thus the cycles of $\lambda\overline d$ are the cycles of
$\overline d$, and its boundaries are $\lambda\mathsf B^q$.
Lemma~\ref{lem:vol2-coinvariant-cohomology} gives the groups.
Finally $\Gamma_\lambda$ is closed, since its two nonzero
differentials are $\lambda dy\xi_1$ and $-\lambda dy\xi_1$,
and its image in the quotient is $[J]$.
\end{proof}

This is a linear quotient calculation. Multiplication by $y$ does
not descend to $\mathsf W$, and no multiplicative quotient of
$\mathcal C_\lambda$ has been asserted.

\subsection{An equivalence of arrows}

Over the level ring use the actual full operator modules
\[
 \mathcal N_R=\operatorname{End}_{R[y]}(R[y]\otimes_kV),
 \qquad
 \mathcal B_R=\operatorname{End}_{D_y\otimes R}(R[y]\otimes_kV).
\]
These permit different parameter degrees in different columns;
they are not replaced by tensoring $N$ or $B_0$ with $R$.
The tower and column recurrences work directly over $R$, giving
\[
 G(\mathcal N_R)\simeq
 \mathcal B_R\oplus(\mathsf W^0\otimes_kR)[-1].
\]
Let $\rho:\mathcal C_\lambda\to G(\mathcal N_R)$ be ghost deletion
followed by multiplication. Define its fibre by
\begin{equation}\label{eq:vol2-integral-fibre-convention}
 \mathcal F_\lambda^n=\mathcal C_\lambda^n\oplus
                  G(\mathcal N_R)^{n-1},\qquad
 d(q,A)=(\mathsf D_\lambda q,\rho q-d_GA).
\end{equation}
Its second coordinate is the retained comparison path.

The contractions of source and target alone would not identify
this fibre. We also need their comparison on the arrow $\rho$.
Vacuum evaluation splits
$\mathcal N_R=\sigma(\mathcal V_R)\oplus\mathcal N_{0,R}$,
with $\sigma(v)=vP_0$ and $\mathcal N_{0,R}1=0$.
On the first summand, evaluation of $\rho$ is exactly ghost
deletion, so the translation quotients commute.

On the second summand, column recursion contracts positive root
degree. Its remaining horizontal component is homotopic to
$a\mapsto a(1-P_0)$ on constants and zero on the positive
quotient. Here is the homotopy. On $dy\,v$ with no ghosts, take
the horizontal operator whose vacuum value is zero and whose
primitive-column value on $u$ is $H_{\mathcal V}(v)u$; extend by
$T^nu\mapsto\nabla^n(H_{\mathcal V}(v)u)$. On positive ghost
inputs it is zero. The identity
\[
 H_{\mathcal V}\nabla f=f-\epsilon(f)(0)
\]
shows that its composite with $\mathsf D_\lambda$ is exactly
the horizontal component of multiplication by $f$ minus
$\epsilon(f)(0)(1-P_0)$. Here $\epsilon$ extracts the constant
state. The first summand contributes $aP_0$. Thus the arrow is
equivalent, with its homotopy supplied, to
\begin{equation}\label{eq:vol2-integral-reduced-arrow}
 R\oplus(\mathsf W_R,\lambda\overline d)[-1]
   \longrightarrow\mathcal B_R\oplus\mathsf W_R^0[-1],
 \qquad (a,w)\longmapsto(a\operatorname{id},\pi_0w),
\end{equation}
where $\mathsf W_R=\mathsf W\otimes_kR$ and $\pi_0$ is the
brutal degree-zero quotient.

\begin{theorem}\label{thm:vol2-integral-fibre}
There is a constructed equivalence of complexes
\begin{equation}\label{eq:vol2-integral-fibre-model}
 \mathcal F_\lambda\simeq
 (\mathcal B_R/R\operatorname{id})[-1]\oplus
       (\mathsf W_R^{\geq1},\lambda\overline d)[-1],
\end{equation}
where $\mathsf W^{\geq1}$ is the brutal subcomplex. Its cohomology is
\begin{align}\label{eq:vol2-integral-fibre-cohomology}
 H^1\mathcal F_\lambda&=\mathcal B_R/R\operatorname{id},\notag\\
 H^2\mathcal F_\lambda&=\mathsf B^1\otimes_kR,\notag\\
 H^n\mathcal F_\lambda&=\mathsf B^{n-1}\otimes_kR/(\lambda)
                              \quad(n\geq3),
\end{align}
and all other groups vanish.
\end{theorem}
\begin{proof}
In \eqref{eq:vol2-integral-reduced-arrow}, the scalar inclusion
splits by vacuum constant evaluation. The second map is
degreewise onto, so its fibre is its kernel
$\mathsf W_R^{\geq1}[-1]$. Explicitly, the redundant two terms
in its cone are the graph
$w\mapsto(-\lambda\overline dw,w)$.
Projection $(v,u)\mapsto v+\lambda\overline du$ and homotopy
$(v,u)\mapsto u$ cancel that graph without dividing by $\lambda$.
This proves the equivalence. In degree two the incoming
$\lambda\mathsf B^1$ disappeared with the cancelled term,
leaving the free module $\mathsf B^1_R$. In higher degrees the
boundaries remain $\lambda\mathsf B^{n-1}$, giving the formula.
\end{proof}

The exact sequence in degree two is therefore
\begin{equation}\label{eq:vol2-integral-path-extension}
 0\longrightarrow(\mathsf W^0/(k\cdot[J]))\otimes_kR
 \xrightarrow{\lambda\overline d}\mathsf B^1\otimes_kR
 \longrightarrow\mathsf B^1\otimes_kR/(\lambda)
 \longrightarrow0.
\end{equation}
Under the isomorphism
$\overline d:\mathsf W^0/(k\cdot[J])\to\mathsf B^1$, the
first map is multiplication by $\lambda$. A concrete instance
uses $F=J^2/2$, for which $d_0F=x_1\xi_1$. The fibre differential
gives
\[
 d(-dy\,F,0)=(\lambda dy\,x_1\xi_1,-dy\,L_F).
\]
Hence
\begin{equation}\label{eq:vol2-integral-path-witness}
 \boxed{\lambda[(dy\,x_1\xi_1,0)]
                 =[(0,dy\,L_{J^2/2})].}
\end{equation}
After its path is forgotten, the current class on the left is
$\lambda$-torsion. In the fibre its lift is a nonzero free class.
Indeed, $J^2$ is a translation primitive independent of $J$,
and the kernel of $\overline d$ in degree zero is only $k\cdot[J]$.

The integral charge-lifting test follows from the same arrow:
\begin{equation}\label{eq:vol2-integral-lift-test}
 [dy\,A]\text{ lifts through }\rho
 \quad\Longleftrightarrow\quad
 \operatorname{red}_+(A1)\in RJ.
\end{equation}
If the normal form is $a(\lambda)J$, the lift is
$a\Gamma_\lambda$ and the comparison homotopy is
$H_N(A-aL_J)$. No inverse level is used. Once a base lift for
a horizontal source shifted by $s$ has been chosen, its lift
space has homotopy groups $H^{s-j}\mathcal F_\lambda$ in
degrees $j\geq0$. Without a base lift the fibre may be empty.

\subsection{The singular level}

The terms of $\mathsf W_R$ are free over $R$.
The quotient $\mathcal B_R/R\operatorname{id}$ is a direct
summand of the torsion-free module $\mathcal B_R$ and hence is
flat over the PID $R$. Derived specialization is therefore
computed by reducing the displayed fibre complex. At $\lambda=0$
its positive differential vanishes, giving
\begin{equation}\label{eq:vol2-zero-level-fibre}
 H^{q+1}(\mathcal F_\lambda\otimes_R^Lk)=\mathsf W^q
                              \qquad(q\geq1).
\end{equation}
In degree two the new classes occur in
\[
 0\longrightarrow\mathsf B^1\longrightarrow\mathsf W^1
      \xrightarrow{\overline d}\mathsf B^2\longrightarrow0.
\]
They are the Tor contribution of integral $H^3$.

For example, $[\xi_0\xi_1]\ne0$ in $\mathsf W^2$.
A polynomial preimage under $\nabla$ would have root degree
zero, by injectivity of $T$ on $\mathcal Q_+$. It would then
give a two-ghost translation preimage in energy zero, where
there is no such monomial. Since
$d_0(x_1\xi_0)=-\xi_0\xi_1$, this is a nonzero element of
$\mathsf B^2$. Thus $dy\,x_1\xi_0$ becomes a degree-two
cocycle at zero level without being an integral cocycle.
The same torsion prevents an $R$-linear deformation retraction
of this fibre complex onto its cohomology: its nonzero
$\lambda$-torsion could not have a cochain section into the
torsion-free terms. This excludes that linear retraction,
not every possible derived replacement.

\subsection{The strict mode-preserving gauge groupoid}

Complete only a new deformation parameter $g$. On
$\mathcal V_R[[g]]$ fix
\[
 J_{-n}=L_{x_n},\qquad J_n=\lambda n\partial_{x_n}\quad(n>0),
 \qquad J_0=q(g)\operatorname{id},\quad q(g)\in gR[[g]].
\]
Seek connections $\nabla+A(g)$ with
$A\in g\mathcal N_R[[g]]$ satisfying
$[\nabla+A,J_n]=-nJ_{n-1}$ for every integer $n$.
A mode-preserving gauge is an element of
$1+g\mathcal N_R[[g]]$ commuting with every nonzero mode.

\begin{theorem}\label{thm:vol2-integral-gauge}
Such a connection exists exactly when $q(g)=\lambda a(g)$
for $a(g)\in gR[[g]]$. Every one is
\begin{equation}\label{eq:vol2-integral-gauge-normal-form}
 \partial_y+T+a(g)L_J+c(y,g)\operatorname{id},
                 \qquad c(y,g)\in gR[y][[g]].
\end{equation}
It is mode-preserving gauge equivalent to $\partial_y+T+a(g)L_J$.
The stabilizer of that normal form is $1+gR[[g]]$.
Thus the strict formal groupoid is the disjoint union, indexed
by $a(g)$, of the classifying groupoids of these scalar units.
\end{theorem}
\begin{proof}
The negative-mode equations force $A$ to commute with all
multiplications by $x_n$. These generate the state module
from $1$, so $A=L_f$. The equations for $n\geq2$ give
$\lambda\partial_{x_n}f=0$. Torsion-freeness makes $f$
independent of those variables. The $n=1$ equation is
$\lambda\partial_{x_1}f=q$. Polynomial coefficient comparison
forces $q=\lambda a$ and $f=aJ+c$. Conversely this formula
satisfies every mode relation.

A gauge commuting with the negative modes is multiplication
by a formal unit. Commutation with the positive modes and
torsion-freeness make it independent of the states. The unit
\[
 U=\exp\left(-\int_0^yc(t,g)\,dt\right)
\]
has polynomial coefficients in every $g$-degree and satisfies
$(\nabla+aL_J+c)U=U(\nabla+aL_J)$.
Between two normal forms a state-independent gauge cannot
change $a$, and its derivative must be zero. These are exactly
the stated scalar stabilizers.
\end{proof}

This groupoid fixes all nonzero modes. It is not the full
derived deformation groupoid, nor the quotient by gauges that
change those modes. At zero level the positive modes vanish;
then arbitrary multiplication corrections are permitted when
$q=0$. Most are not specializations of
\eqref{eq:vol2-integral-gauge-normal-form}.

Root completion also changes this problem. The series
$g_0=\sum_{n\geq0}(-1)^ny^{n+1}T^nJ/(n+1)!$ satisfies
$\nabla g_0=J$, so the integral charge image $dy\,L_J$ becomes
exact. In the completed $J$-tower the section
$1\mapsto e^{-yT}J=J-Tg_0$ splits the quotient by state
translation. This is the completion of the polynomial
extension $0\to P\xrightarrow{R_\partial}P\to O\to0$.
The polynomial projective model need not remain projective in
the continuous category. Its nonzero Massey coset and its
integral path relation therefore cannot be transported by
an undeclared equivalence across this change of coefficients.

\section{Finite categories of interfaces}
\label{sec:vol2-finite-interface-centres}

An invertible bimodule composes by derived tensor product, as in
\eqref{eq:vol2-interface-composition}. A collection of interfaces
with specified junction maps has more data: the centre must act
compatibly with every composition of those maps. For two objects,
Proposition~\ref{prop:vol2-dg-triangular} expresses this as the fibre
of the difference between the two endpoint characteristic actions
on the joining bimodule. With more objects, the longer composable
strings contribute further terms.

Let $R$ be a commutative $k$-algebra and let $\mathcal A$ be a dg
category with ordered objects $0,\ldots,n-1$, scalar endomorphisms
$\mathcal A(i,i)=R$, and $\mathcal A(i,j)=0$ for $i>j$.
Suppose each forward complex $M_{ij}$ is bounded and finite
projective over $R$, and all dg compositions are specified and
associative. Let $A$ be the sum algebra, $S=\bigoplus_iRe_i$,
and $\overline A=\ker(A\to S)$. The normalized relative bimodule
bar is
\begin{equation}\label{eq:vol2-directed-bar}
 P_\ell=A\otimes_S(\overline A[1])^{\otimes_S\ell}\otimes_SA,
                     \qquad0\leq\ell\leq n-1.
\end{equation}
For an unsuspended tensor of internal degree $q$ and bar length
$\ell$, the total degree is $q-\ell$. One sign convention is
$d_{\mathrm{int}}+(-1)^q\sum_{j=0}^{\ell}(-1)^jd_j$, where
$d_j$ multiplies adjacent factors, including the two endpoint
actions. Suspension translates this to
\eqref{eq:vol2-directed-bar}.

\begin{theorem}\label{thm:vol2-directed-bar}
The total complex $P$ is a finite homotopically projective
$A$-bimodule resolution of $A$. Thus the relative derived centre is
\[
 Z_R(A)=R\operatorname{Hom}_{A\otimes_RA^{\mathrm{op}}}(A,A)
          \simeq\operatorname{Hom}_{A^e}(P,A).
\]
Its evaluation gives, for every complex $F$ of $R$-modules,
\begin{equation}\label{eq:vol2-directed-evaluation}
 R\operatorname{Hom}_R(F,Z_R(A))
 \simeq R\operatorname{Hom}_{A^e}(F\otimes_R^LA,A).
\end{equation}
For every derived commutative scalar extension $R\to R'$,
\begin{equation}\label{eq:vol2-directed-base-change}
 Z_R(A)\otimes_R^LR'
       \simeq Z_{R'}(A\otimes_R^LR'),
\end{equation}
compatibly with evaluation, cup products and Hochschild insertions.
\end{theorem}
\begin{proof}
A nonzero string of $\ell$ forward arrows visits $\ell+1$
strictly increasing objects, so $\ell\leq n-1$.
The unit-insertion contraction of the augmented unnormalized
bar proves its exactness after forgetting an outer action.
The degenerate subcomplex splits off, leaving the normalized
resolution.

Each remaining term is induced from a tensor of bounded finite
projective coefficient complexes and two representable
idempotent modules. Hom into an acyclic bimodule reduces to
Hom from that coefficient complex into an acyclic idempotent
component, hence is acyclic. The finite totalization is therefore
homotopically projective. Tensor--Hom adjunction applied to
$F\otimes_RP$ proves \eqref{eq:vol2-directed-evaluation}.

Every cochain coefficient is a Hom between bounded finite
projective $R$-complexes. These Homs and the finite sums commute
with derived scalar extension. The bar differential, cuts and
insertions use the same composition maps after extension.
This proves \eqref{eq:vol2-directed-base-change} with its operations.
\end{proof}

Strictly increasing object indices made the bar finite.
With arbitrary dg endpoint algebras the repeated corner inputs
of Proposition~\ref{prop:vol2-dg-triangular} instead require
product totalization. We now use the finite case to see what
the operations and the parameter fibres retain.

\section{Koszul parameters and contraction brackets}
\label{sec:vol2-koszul-interface}

Let $r\geq1$, $R=k[h_1,\ldots,h_r]$, $I=(h_1,\ldots,h_r)$,
and $S_0=R/I$. Let $U_0=k^r$ have degree-one basis
$u_1,\ldots,u_r$, and write $L_i$ for exterior multiplication
by $u_i$. Take the arrow complex
\begin{equation}\label{eq:vol2-koszul-arrow}
 \begin{gathered}
 M=R\otimes_k\Lambda U_0,\qquad d_M=\sum_ih_iL_i,\\
 T(M)=Re_0\oplus Re_1\oplus M,\qquad e_1Me_0=M,\quad M^2=0.
 \end{gathered}
\end{equation}
The differential is a differential of the arrow module; it is
not a derivation of the auxiliary exterior multiplication.
Anticommutation gives $d_M^2=0$, and the Koszul complex has
cohomology $S_0$ in degree $r$.

Put $E=\operatorname{End}_R(M)$ with
$\delta f=d_Mf-(-1)^{|f|}fd_M$. Reverse the one-arrow coordinate
in Theorem~\ref{thm:vol2-triangular-cone}. The centre then has
the finite model
\begin{equation}\label{eq:vol2-koszul-centre}
 Z^n=(R^2)^n\oplus E^{n-1},\qquad
 d(a,b;f)=(0,0;(b-a)\operatorname{id}-\delta f).
\end{equation}
The endpoint terms occur only in degree zero. Use the
source-first opposite-cup convention
\begin{equation}\label{eq:vol2-koszul-cup}
 (a,b;f)(a',b';f')=(aa',bb';fb'+af').
\end{equation}
In this convention, for $\sigma f=(0,0;f)$ of degree $|f|+1$,
the insertion bracket is
\begin{equation}\label{eq:vol2-koszul-bracket}
 [\sigma f,\sigma g]
       =\sigma\bigl(fg-(-1)^{|f||g|}gf\bigr).
\end{equation}
Endpoint scalars have zero bracket with the mixed terms. The
bracket has degree $-1$.

These operations follow directly from the one-arrow bar. Its
two cuts give \eqref{eq:vol2-koszul-cup}, while insertion into
the only arrow input is composition in $E$. Two mixed cup
factors would require two composable forward arrows, which do
not exist. The differential identity on endpoints is
$bb'-aa'=(b-a)b'+a(b'-a')$; the other identities reduce to
$R$-linearity of $\delta$. On $Z[1]$ the mixed differential
is $\delta$, so the bracket differential and Jacobi identities
are the graded-commutator identities. This fixes the signs
of the comparison as well as its underlying complex.

\subsection{The integral ring and its bracket}

Let $\iota_i$ contract by the dual of $u_i$. The relations are
\begin{equation}\label{eq:vol2-koszul-clifford}
 \{L_i,L_j\}=0,\qquad \{\iota_i,\iota_j\}=0,
 \qquad \{\iota_i,L_j\}=\delta_{ij}\operatorname{id}.
\end{equation}
The operators $L_A\iota_B$, with increasing subsets $A,B$,
form a basis of $E$: for one generator $1,L,\iota,L\iota$
span all matrices, and graded tensor products give the general
case. Moreover $\delta L_i=0$ and
$\delta\iota_i=h_i\operatorname{id}$. At fixed creation
monomial $L_A$, the contraction variables are the Koszul
resolution of $(h_i)$. Therefore
\[
 H^q(E)=S_0\otimes_k\Lambda^qU_0\quad(0\leq q\leq r),
 \qquad H^q(E)=0\text{ otherwise}.
\]

\begin{theorem}\label{thm:vol2-koszul-centre-ring}
The relative centre of \eqref{eq:vol2-koszul-arrow} has
\begin{align*}
 \operatorname{HH}_R^0(T)&=R\times_{S_0}R,&
 \operatorname{HH}_R^1(T)&=0,\\
 \operatorname{HH}_R^{q+1}(T)&=S_0\otimes_k\Lambda^qU_0
                         &&(1\leq q\leq r),
\end{align*}
and no other cohomology. Its positive-degree cup ideal is
square-zero; a pair in degree zero acts through its common
residue in $S_0$. For $t_i=[(0,h_i;\iota_i)]$,
\begin{equation}\label{eq:vol2-koszul-centre-presentation}
 \operatorname{HH}_R^0(T)
 =R[t_1,\ldots,t_r]/(h_it_j-h_jt_i,\ t_it_j-h_it_j)_{i,j}.
\end{equation}
\end{theorem}
\begin{proof}
The map $R^2\to H^0(E)=S_0$ is $(a,b)\mapsto b-a$.
It is onto with kernel $R\times_{S_0}R$, so the fibre exact
sequence gives the groups. The mixed cup product is zero;
the endpoint residue gives its scalar action.

The representative of $t_i$ is closed because
$\delta\iota_i=h_i$. It maps to $(0,h_i)$ in the fibre product.
Every pair is a diagonal scalar plus $(0,c)$ with $c\in I$,
so these classes generate. The quadratic relations reduce
each polynomial to $a+\sum_ib_it_i$. It maps to zero exactly
when $a=0$ and $\sum_ib_ih_i=0$. The first Koszul syzygies
$h_it_j-h_jt_i$ generate this kernel, proving the presentation.
\end{proof}

For $A=\{a_1<\cdots<a_q\}$ set
$\alpha_A=[\sigma L_{a_1}\cdots L_{a_q}]$, of degree $q+1$.
Then for $q\geq2$,
\begin{equation}\label{eq:vol2-koszul-contraction-brackets}
 [t_i,\alpha_A]=
 \begin{cases}
 (-1)^{j-1}\alpha_{A\setminus\{i\}},&i=a_j,\\
 0,&i\notin A.
 \end{cases}
\end{equation}
Indeed, the graded commutator of $\iota_i$ with the creation
product is precisely this contraction. When $q=1$ the result
is the scalar identity, which is the boundary of an endpoint
difference in \eqref{eq:vol2-koszul-centre}, so the bracket
is zero. All brackets between two positive classes or two
degree-zero classes vanish; the generators supercommute in
their respective creation or contraction factors.

For $r=2$, with $\omega=\alpha_{\{1,2\}}$, this gives
\[
 [t_1,\omega]=\alpha_{\{2\}},\qquad
 [t_2,\omega]=-\alpha_{\{1\}}.
\]
The same $t_i$ annihilate every positive class under cup
product. Thus the bracket detects a composition invisible
to that product. The proof of the integral calculation also
works for any regular sequence $(h_i)$ in a commutative
$\mathbb Q$-algebra: regularity is used in the Koszul
cohomology and in its first syzygies.

\subsection{The full singular fibre}

The finite bar commutes with derived scalar extension by
Theorem~\ref{thm:vol2-directed-bar}. If some $h_i$ is invertible,
$h_i^{-1}\iota_i$ contracts the arrow $M$ and the centre becomes
the two scalar endpoint sectors. At the common zero, put
$V_0=\Lambda U_0$. The arrow differential vanishes, and the
only centre differential sends the endpoint difference to
$\operatorname{id}\in E^0$, in centre degree one. Hence
\begin{equation}\label{eq:vol2-koszul-singular-fibre}
 \operatorname{HH}_k^*(T\otimes_Rk)
   =k\ltimes(\operatorname{End}_k(V_0)/k\operatorname{id})[-1].
\end{equation}
Its displayed ideal has square-zero cup product and shifted
graded-commutator bracket. The unit has zero bracket.
Since
\[
 \dim_k\operatorname{End}^p(\Lambda U_0)
    =\sum_j\binom rj\binom r{j+p}=\binom{2r}{r+p},
\]
the dimensions are
\begin{equation}\label{eq:vol2-koszul-fibre-dimensions}
 \dim_k\operatorname{HH}^n_k(T\otimes_Rk)=
 \begin{cases}
 1+\binom{2r}{r-1},&n=0,\\
 \binom{2r}{r}-1,&n=1,\\
 \binom{2r}{r+n-1},&n\ne0,1.
 \end{cases}
\end{equation}
Out-of-range binomial coefficients are zero. The total dimension
is $4^r$. For $r=2$, degrees $-1,0,1,2,3$ have dimensions
$1,5,5,4,1$. In contrast, ordinary specialization of integral
degree-zero cohomology has basis $1,t_1,t_2$ and dimension three.
Specializing the complex retains the missing extension data.

There is no $R$-linear deformation retraction of
\eqref{eq:vol2-koszul-centre} onto its integral cohomology.
The nonzero $\operatorname{HH}_R^2(T)$ is annihilated by $h_1$,
whereas multiplication by $h_1$ is injective on each free
cochain term. A cochain section of this torsion group must be
zero and cannot have a left inverse. This argument concerns
a linear retraction over $R$; it is not a classification of
all formality zigzags.

\subsection{A second-order obstruction}

At $r=2$ and zero parameters, use the basis
$v_0=1,v_1=u_1,v_2=u_2,v_3=u_1u_2$ of $V_0$. Define the
degree-one endomorphism $f$ by
\[
 f(v_0)=v_1,\qquad f(v_1)=v_3,\qquad f(v_2)=f(v_3)=0.
\]
Over $k[\epsilon]/(\epsilon^2)$, the arrow differential
$\epsilon f$ defines a first-order deformation of $T_0$.
But $f^2(v_0)=v_3$, and
\begin{equation}\label{eq:vol2-koszul-obstruction}
 \tfrac12[\sigma f,\sigma f]=\sigma f^2\ne0
                         \quad\text{in }\operatorname{HH}^3_k(T_0).
\end{equation}
It is nonzero in \eqref{eq:vol2-koszul-singular-fibre} because
an internal degree-two map is not a scalar identity.
Thus the second-order Hochschild equation
$d_{\mathrm{Hoch}}g+\tfrac12[\sigma f,\sigma f]=0$ has no
solution. This excludes even curvature-free $A_\infty$
corrections on the same graded module with the same
first-order class, although the cup square of $\sigma f$
is zero. By contrast, the original Koszul directions
$xL_1+yL_2$ square to zero and integrate over $k[x,y]$.

\section{Four composable objects}
\label{sec:vol2-four-interfaces}

The preceding category had one arrow direction and no composable
pair of nonidentity arrows. Four ordered objects provide the
first three-arrow string. Its two composition orders give an
associativity equation, and the centre retains that equation
in its own differential.

Use chronological products in this section: $a_{ij}a_{jk}$ means
first $i\to j$, then $j\to k$, so $e_ia_{ij}=a_{ij}=a_{ij}e_j$.
Take objects $0<1<2<3$, scalar identities, and a degree-zero
basis arrow $a_{ij}$ for each $i<j$. Set
\begin{equation}\label{eq:vol2-four-products}
 a_{01}a_{12}=pa_{02},\quad a_{12}a_{23}=qa_{13},\quad
 a_{02}a_{23}=ra_{03},\quad a_{01}a_{13}=sa_{03}.
\end{equation}
All other products are forced by identities or are zero.
The only associator of nonidentity arrows is
\[
 (a_{01}a_{12})a_{23}-a_{01}(a_{12}a_{23})
       =Fa_{03},\qquad F=pr-qs.
\]
Thus the associative family is the rank-ten algebra $A_B$ over
$B=k[p,q,r,s]/(pr-qs)$.

\subsection{The centre differential and its operations}

Order the edges as $01,02,03,12,13,23$ and the triangles as
$012,123,023,013$. A normalized cochain assigns a scalar to
each increasing string, multiplying its endpoint arrow.
The finite bar gives
\begin{equation}\label{eq:vol2-four-centre}
 Z_B(A_B)\simeq
 [B^4\xrightarrow{d_0}B^6\xrightarrow{d_1}B^4\xrightarrow{d_2}B]
\end{equation}
in degrees zero through three, with
\begin{align}\label{eq:vol2-four-differentials}
 (d_0\lambda)_{ij}&=\lambda_j-\lambda_i,\notag\\
 d_1f&=\bigl(p(f_{01}+f_{12}-f_{02}),
             q(f_{12}+f_{23}-f_{13}),\notag\\
 &\hspace{20mm}r(f_{02}+f_{23}-f_{03}),
             s(f_{01}+f_{13}-f_{03})\bigr),\notag\\
 d_2(\alpha,\beta,\gamma,\delta)&=-r\alpha+s\beta-p\gamma+q\delta.
\end{align}
Indeed, in arity zero the two Hochschild endpoint terms give
$d_0$. On a triangle the three faces have the same composition
scalar, giving $d_1$. On $0123$ the four terms give $d_2$.
Before imposing the relation, direct substitution gives
\begin{equation}\label{eq:vol2-four-square}
 d_1d_0=0,\qquad
 d_2d_1f=-F(f_{01}+f_{12}+f_{23}-f_{03}).
\end{equation}
The square-zero condition remembers the original associator.

Write $m_{012}=p,m_{123}=q,m_{023}=r,m_{013}=s$.
For positive arities $n,m$, the cup product is
\begin{equation}\label{eq:vol2-four-cup}
 (f\smile g)_{i_0\cdots i_{n+m}}
   =m_{i_0i_ni_{n+m}}f_{i_0\cdots i_n}g_{i_n\cdots i_{n+m}}.
\end{equation}
An arity-zero factor acts by the appropriate first or last
endpoint scalar. A product with no increasing string is zero.
The bar cut proves this formula; two successive cuts prove
associativity using $pr=qs$, and the faces give the Leibniz
identity. Its unit is $(1,1,1,1)$.

For insertion, use
$f\circ g=\sum_{j=0}^{n-1}(-1)^{j(m-1)}f\circ_jg$ and
$[f,g]=f\circ g-(-1)^{(n-1)(m-1)}g\circ f$.
For unary $f$ and binary $g$ this gives
\[
 [f,g]_{ijk}=(f_{ik}-f_{ij}-f_{jk})g_{ijk}.
\]
For binary cochains it gives
\begin{equation}\label{eq:vol2-four-bracket}
 [f,g]_{0123}=f_{023}g_{012}-f_{013}g_{123}
                   +g_{023}f_{012}-g_{013}f_{123}.
\end{equation}
These are the two insertion slots on the unique three-arrow
string. With the Hochschild signs of
\eqref{eq:vol2-four-differentials}, a positive arity-$n$
cochain has $df=(-1)^{n-1}[\mu,f]$, where $\mu$ is the binary
product. Rephasing arities gives the dg Lie convention
$[\mu,-]$. In particular $[\mu,\mu]/2=Fa_{03}$.

\subsection{The parameter fibres and the integral module}

Let $B\to K$ be a map to a field, and let $c$ count the nonzero
parameters among $p,q,r,s$. Theorem~\ref{thm:vol2-directed-bar}
makes the specialized complex the full centre over $K$.
Its cohomology dimensions are
\begin{equation}\label{eq:vol2-four-fibres}
 \begin{array}{c|rrrr}
 c&\operatorname{HH}^0&\operatorname{HH}^1&
          \operatorname{HH}^2&\operatorname{HH}^3\\\hline
 0&1&3&4&1\\
 1&1&2&2&0\\
 2&1&1&1&0\\
 4&1&0&0&0
 \end{array}
\end{equation}
The proof is a rank calculation. The incidence matrix $d_0$
has rank three. The four unweighted triangle rows of $d_1$
have one relation with all four coefficients nonzero, and
any three are independent. Thus $\operatorname{rank}d_1$
is $\min(c,3)$. The last row has rank one off the origin
and zero at the origin. Subtracting adjacent ranks gives
the table. The equation $pr=qs$ excludes $c=3$ and excludes
the supports $\{p,r\}$ and $\{q,s\}$ when $c=2$.

Integrally $B$ is a domain: $pr-qs$ is primitive and
irreducible as a polynomial in $r$ over $k[p,q,s]$.
The answer retains a module supported at the origin:
\begin{align}\label{eq:vol2-four-integral}
 \operatorname{HH}_B^0(A_B)&=B,\qquad
 \operatorname{HH}_B^1(A_B)=0,\notag\\
 \operatorname{HH}_B^3(A_B)&=B/(p,q,r,s),\notag\\
 \operatorname{HH}_B^2(A_B)&=
   \ker(-r,s,-p,q)/\operatorname{im}L.
\end{align}
Here $L(x,y,z)=(-px,-qy,r(x-z),s(y-z))$.
For the degree-one assertion, nonzero parameters in the domain
force all four triangle additivity equations. Setting
$\lambda_0=0$ and $\lambda_i=f_{0i}$ gives
$f_{ij}=\lambda_j-\lambda_i$. More generally, modulo a
vertex coboundary an edge cochain has a unique representative
with $f_{01}=f_{12}=f_{23}=0$. Its other coordinates
$(f_{02},f_{13},f_{03})=(x,y,z)$ give the displayed image
$L$. The cokernel of $d_2$ gives degree three.

There is no $B$-linear deformation retraction onto these
cohomology groups. Its nonzero degree-three group is killed
by $p$, and therefore cannot have a section into the free
degree-three cochain term. Scalar extension must act on
\eqref{eq:vol2-four-centre}, not merely on
\eqref{eq:vol2-four-integral}.

At the origin, let the binary cochain $f$ have
$f_{012}=f_{023}=1$ and $f_{123}=f_{013}=0$.
The first-order product $\epsilon f$ is associative modulo
$\epsilon^2$, but
\[
 \tfrac12[f,f]_{0123}=1\in\operatorname{HH}^3=k.
\]
The positive-arity differential is zero there, so no
second-order correction removes this obstruction. Equivalently,
$p=r=\epsilon$, $q=s=0$ gives $F=\epsilon^2$.
A ternary operation cannot help on the unchanged degree-zero
arrow module over an ordinary ring: it would have degree $-1$
and no target. This is the deformation problem with the
objects, arrow bases and units fixed.

\subsection{A primitive and a finite coherent resolution}

Introduce a coefficient of degree $-1$ with the required
differential:
\begin{equation}\label{eq:vol2-four-coherent-base}
 R_{\mathrm{coh}}=k[p,q,r,s]\otimes\Lambda(\chi),
                  \qquad d\chi=pr-qs.
\end{equation}
Its cohomology is $B$ in degree zero, since $F$ is a
non-zero-divisor. A map to a derived test algebra nevertheless
retains the image of $\chi$ as a specified primitive.
On the rank-ten arrow module, the products
\eqref{eq:vol2-four-products} and
\begin{equation}\label{eq:vol2-four-ternary}
 m_3(a_{01},a_{12},a_{23})=-\chi a_{03}
\end{equation}
define a strictly unital $A_\infty$ category, with convention
$dm_3=-\operatorname{Assoc}(m_2)$ on closed degree-zero
arrows. There is no increasing string of four nonidentity
inputs, so all higher normalized identities vanish.

The binary homotopies and their compatibility can also be
retained in a strict dg algebra. Let $\mathcal Q$ be the
free directed path algebra over $R_{\mathrm{coh}}$ on
\[
 X_{ij}:i\to j\ (|X_{ij}|=0),\qquad
 H_{ijk}:i\to k\ (|H_{ijk}|=-1),\qquad
 K_{03}:0\to3\ (|K_{03}|=-2),
\]
for increasing indices. Extend the coefficient differential
by the graded Leibniz rule and
\begin{align}\label{eq:vol2-four-resolution}
 dX_{ij}&=0,\notag\\
 dH_{ijk}&=X_{ij}X_{jk}-m_{ijk}X_{ik},\notag\\
 dK_{03}&=H_{012}X_{23}-X_{01}H_{123}
                 +pH_{023}-qH_{013}+\chi X_{03}.
\end{align}
Odd coefficients are moved to the left with their Koszul sign.
In $d^2K_{03}$ the first four terms cancel all paths except
$-FX_{03}$; the last gives $FX_{03}$. Thus $d^2=0$.
The free basis has fifteen degree-zero paths (four identities,
six arrows, four length-two paths and one length-three path),
six degree-$-1$ paths (the four $H$ and their two composable
extensions), and the single $K_{03}$. Its rank is $22$.

Let $V_{\mathrm{ar}}$ be the submodule spanned by the identities
and the six $X_{ij}$, with its coefficient differential.
Let $i$ include it, and let $\pi$ be the identity there and
zero on the $H_{ijk},K_{03}$. Specify its remaining values by
\begin{align}\label{eq:vol2-four-projection}
 \pi(X_{ij}X_{jk})&=m_{ijk}X_{ik},\notag\\
 \pi(X_{01}X_{12}X_{23})&=prX_{03},&
 \pi(X_{01}H_{123})&=\chi X_{03},&
 \pi(H_{012}X_{23})&=0.
\end{align}
The degree-$-1$ homotopy is
\begin{align}\label{eq:vol2-four-homotopy}
 h(X_{ij}X_{jk})&=H_{ijk},\notag\\
 h(X_{01}X_{12}X_{23})&=H_{012}X_{23}+pH_{023},&
 h(X_{01}H_{123})&=-K_{03},
\end{align}
and is zero on all other basis paths. Extend $\pi$ linearly
and $h$ graded-linearly, so $h(\chi a)=-\chi h(a)$.

\begin{theorem}\label{thm:vol2-four-contraction}
Over all of $R_{\mathrm{coh}}$ these maps satisfy
\[
 \pi i=1,\qquad dh+hd=1-i\pi,\qquad
 h^2=hi=\pi h=0,\qquad d\pi=\pi d.
\]
The transferred product on $V_{\mathrm{ar}}$ is
\eqref{eq:vol2-four-products}, with the ternary value
\eqref{eq:vol2-four-ternary}; every normalized higher
operation vanishes. No parameter is inverted.
\end{theorem}
\begin{proof}
The intervals $0\to2$ and $1\to3$ contain the pairs
$H_{012}\mapsto X_{01}X_{12}-pX_{02}$ and
$H_{123}\mapsto X_{12}X_{23}-qX_{13}$.
On $0\to3$ put
\begin{align*}
 v_1&=X_{02}X_{23}-rX_{03},& h_1&=H_{023},\\
 v_2&=X_{01}X_{13}-sX_{03},& h_2&=H_{013},\\
 v_3&=X_{01}X_{12}X_{23}-pX_{02}X_{23},& h_3&=H_{012}X_{23},\\
 \ell&=H_{012}X_{23}-X_{01}H_{123}
                   +pH_{023}-qH_{013}+\chi X_{03}.
\end{align*}
Together with $K_{03},X_{03}$ these are another free basis,
obtained by a triangular change with constant invertible
pivots. Its differential has $dh_j=v_j$, $dK_{03}=\ell$,
$d\ell=0$, leaving $X_{03}$. Contracting these four pairs
and returning to the original basis gives
\eqref{eq:vol2-four-projection}--\eqref{eq:vol2-four-homotopy}.
In particular $\pi(\ell)=0$ forces the nonzero value
$\pi(X_{01}H_{123})=\chi X_{03}$. Graded linearity of $h$
cancels the coefficient-differential terms.

The transferred binary operation is $\pi\mu(i\otimes i)$.
Its two ternary trees give
\[
 \pi\bigl(h(X_{01}X_{12})X_{23}
                   -X_{01}h(X_{12}X_{23})\bigr)=-\chi X_{03}.
\]
There are no longer composable nonidentity strings, and the
side conditions kill higher trees with an identity input.
This proves the strictly unital structure.
\end{proof}

After an ordinary specialization $F=\chi=0$, the augmentation
$X_{ij}\mapsto a_{ij}$, $H_{ijk},K_{03}\mapsto0$ is a strict
dg algebra quasi-isomorphism, by this contraction. Over
$R_{\mathrm{coh}}$ the projection is a cochain map to the
$A_\infty$ category, not a strict algebra map.
The primitive has not removed the ordinary obstruction:
a dg map to $k[\epsilon]/(\epsilon^3)$ with $p=r=\epsilon$
and $q=s=0$ would need a degree-$-1$ primitive of
$\epsilon^2$, which that ring does not contain.

The finite categories in this chapter specify algebraic
compositions of interfaces. Their object labels alone supply
no field-theoretic realization. Such a realization must
identify the arrows, junctions, collision maps and local
vacuum operations. Likewise, a native centre comparison must
carry the full current hierarchy, the polynomial Massey
coset and the integral path relation, together with derived
specialization. These are tests on an actual evaluation map;
they do not construct a common recipient for arbitrary
geometric boundary actions.
